% 03-related.tex -- draft. Re-read Meshlib and Cvetkovic & Klimovic before submission.
\section{Related Work}
\label{sec:related}

\paragraph{Component-level analysis of sidecars}
MeshInsight~\cite{zhu2023meshinsight} is the work closest in method to ours.
It profiles the Istio sidecar (Istio 1.13, Envoy 1.21) on Online Boutique and
Hotel Reservation and attributes its overhead to seven components, finding that
protocol parsing dominates. It reports latency increases of 27--269\% and CPU
increases of 42--163\%, depending on configuration, and proposes a linear
compositional model of overhead. The authors note that the model does not
predict tail latency, and that extending the analysis to other meshes is future
work. We extend the analysis to ztunnel, the Ambient waypoint,
\texttt{linkerd2-proxy}, and Cilium's eBPF datapath, use open-loop load, add
energy, and model the tail through a queueing term.

\paragraph{Comparisons of sidecar and sidecar-less meshes}
Bremler-Barr et al.~\cite{bremlerbarr2024mtls} compare Istio sidecar, Istio
Ambient (ztunnel only), Linkerd, and Cilium, focusing on the cost of mTLS. Their
workload is a synthetic server with a fixed 200\,ms delay driven by Fortio, and
they do not exercise L7 processing, waypoints, or realistic call graphs. They
also observe that Cilium does not encrypt traffic between pods on the same node,
a caveat that applies to our setup as well: under our placement, 2.8 of the
3.8 internal calls per Hotel Reservation request, and 0.6 of the 11 per Online
Boutique request, stay on one node and are therefore not encrypted by Cilium's
WireGuard. Our study uses multi-service
applications, includes the waypoint, and reports energy and a predictive model.

\paragraph{Reducing mesh overhead}
A separate line of work proposes new data-plane designs rather than measuring
existing ones. Meshlib~\cite{mostafaei2026meshlib} enforces mesh policy inside
the application process for sidecar-less meshes, and Cvetkovi\'c and
Klimovic~\cite{cvetkovic2025lightweight} design a lightweight sidecar-based mesh
for serverless platforms, where Ambient is not compatible with the platform.
\todo{re-read both papers and summarize their measured baselines in one
sentence each} Our measurements are complementary: they quantify the cost these
designs aim to remove, on the architectures in production today.

\paragraph{Microservice benchmarks and load generation}
We use Online Boutique~\cite{onlineboutique} and Hotel Reservation from
DeathStarBench~\cite{gan2019deathstarbench}, the same applications as
MeshInsight, so that results can be compared directly. We generate open-loop
load with k6~\cite{k6} at a constant arrival rate, avoiding the
coordinated-omission problem of closed-loop generators, which underestimates
tail latency~\cite{tene2015latency}. Energy is measured with Intel
RAPL~\cite{khan2018rapl}.
\todo{add 1--2 references on energy measurement of cloud-native software, if
space allows}
